% Bewertungspaket zum Gesamttest «Lineare Funktionen» — Quelle des PDFs (python3 scripts/build-lp-pdf.py lineare-funktionen).
% Ganze Punkte je Teilschritt; Werte und Fehlerwerte mit python3 nachgerechnet (scripts/lp/lineare-funktionen/zahlen.py).
% Gerüst aus dem Bewertungspaket von LP 3 (gleichungen-formeln): Auftrag an die KI, Folgefehler, (E)-Zeilen.
\documentclass[11pt]{article}
\usepackage{../lp-druck}
\usepackage{lf-gitter}
\renewcommand{\lpname}{Lineare Funktionen}
\renewcommand{\lpteilgebiet}{Funktionen ·}
\renewcommand{\lpfuss}{Bewertungspaket}
\newcolumntype{P}{>{\centering\arraybackslash}p{10mm}}
\newenvironment{raster}{\par\smallskip\noindent\renewcommand{\arraystretch}{1.6}\linespread{1.35}\selectfont\tabularx{\linewidth}{@{}>{\raggedright\arraybackslash}p{38mm}P X@{}}\toprule
  \small\textsc{Teilschritt} & \small\textsc{P} & \small\textsc{Musterlösung}\\\midrule}{\bottomrule\endtabularx\par}
\newcommand{\fehler}[1]{\par\smallskip{\small\textcolor{lprot}{\bfseries Typische Fehler:} #1}}
\newcommand{\E}{\,{\footnotesize\bfseries\color{lpgruen}(E)}}
\begin{document}
\lpkopf{Bewertungspaket zum Gesamttest}{}

\begin{lpachtung}{Erst lösen, dann öffnen}
Dieses Paket enthält die Musterlösung. Wer es vor dem Lösen liest, misst am Ende nur, wie gut das Abschreiben gelingt.
\end{lpachtung}

\section*{1 · So gehst du vor}
\begin{enumerate}[leftmargin=*,itemsep=2pt]
\item \textbf{Gesamttest lösen} — vollständig, mit Rechenweg, ohne dieses Paket und ohne Taschenrechner, mit Lineal.
\item \textbf{Selbst bewerten} mit Abschnitt~3: Zeile für Zeile vergleichen und die Punkte eintragen. Das geht ohne Hilfsmittel
  und ist der Normalfall.
\item \textbf{Oder bewerten lassen.} Ob du dafür eine KI nutzt und welche, entscheidest du selbst und in eigener Verantwortung —
  je nachdem, was dir aufgrund deines Alters und deiner persönlichen Situation erlaubt ist (Altersgrenzen und
  Nutzungsbedingungen des Anbieters, allenfalls Einverständnis deiner Eltern). Fotografiere oder scanne deine Lösung gut lesbar,
  \textbf{ohne Namen}. Handyfotos tragen oft unsichtbar Standort, Zeit und Gerät mit (Metadaten) — lade lieber einen Scan oder ein
  Bildschirmfoto des Fotos hoch. Lade nur deine Lösung, das Testblatt und dieses PDF hoch, keine persönlichen Daten, und füge den Auftrag aus
  Abschnitt~2 ein.
\item \textbf{Rückmeldung prüfen.} Stimmt, was die KI aus deiner Handschrift gelesen hat? Eine KI kann sich irren —
  im Zweifel gilt das Raster in Abschnitt~3.
\item \textbf{Gezielt wiederholen} nach Abschnitt~4.
\end{enumerate}

\section*{2 · Auftrag an die KI}
\begin{tcolorbox}[colback=lplinie!25,colframe=lplinie,boxrule=0.5pt,arc=1mm,fontupper=\small]
Du bist Korrektorin oder Korrektor für Mathematik auf der Sekundarstufe I in der Schweiz. Im Anhang findest du (1)~das Testblatt mit den Aufgaben, (2)~das Bewertungspaket zum Gesamttest «Lineare Funktionen» mit Musterlösung und Punkteraster und (3)~Fotos meiner handschriftlichen Lösung.\par\medskip
Geh so vor:\par
1. Schreib zu jeder Aufgabe G1 bis G8 zuerst ab, was du in meiner Lösung liest — Rechenweg und Ergebnis; bei den Zeichnungen (G4, G6, G8), durch welche Punkte meine Gerade oder meine Strecken gehen. Was du nicht sicher lesen kannst, markierst du mit [unsicher]. Nicht raten. Achte besonders auf Minuszeichen und Brüche.\par
2. Vergib die Punkte genau nach dem Raster im Bewertungspaket (Abschnitt 3), ganze Punkte je Zeile. Ziehe einen Fehler nur einmal ab: Rechne ich mit einem falschen Zwischenergebnis richtig weiter, gibt es die Folgepunkte. Ausnahme: Zeilen mit \textbf{(E)} gibt es nur für das richtige Ergebnis, nie als Folgepunkt. Die Zeilen «Typische Fehler» sagen, welche Rasterzeile bei diesem Fehler 0 Punkte gibt — sie sind kein zusätzlicher Abzug.\par
3. Andere richtige Lösungswege zählen voll, zum Beispiel eine Wertetabelle statt des Steigungsdreiecks oder bei G7 der andere Punkt zum Einsetzen — \textbf{ausser} die Aufgabe verlangt etwas ausdrücklich: G4 das Nachrechnen der Punktprobe (Ablesen allein genügt nicht), G7 die Probe mit $Q$. Gleichwertig: $y = 1.5x - 2$ und $y = \tfrac{3}{2}x - 2$; $y = -6x + 80$ und $y = 80 - 6x$; Brüche und Dezimalzahlen mit gleichem Wert. Eine Zeichnung zählt, wenn die verlangten Punkte auf etwa ein halbes Häuschen genau getroffen sind.\par
4. Begründe jeden Abzug in einem Satz und nenne die Fehlerart (zum Beispiel «m und b vertauscht», «Bruch der Steigung umgedreht», «Reihenfolge gemischt», «b mit einem y-Wert verwechselt», «nach dem Wechsel neu angefangen»). Schreib die Musterlösung nicht ab — zeig mir, wo mein Fehler liegt.\par
5. Gib zum Schluss eine Tabelle «Aufgabe | Punkte | Begründung», die Gesamtpunktzahl (von 25) und — nach der Selbsteinschätzung in Abschnitt 4 — welches Kapitel des Leitprogramms ich wiederholen soll. Keine Note.\par
6. Fehlt ein Foto oder ist eine Aufgabe unlesbar, sag es, statt sie zu bewerten.
\end{tcolorbox}

\needspace{16\baselineskip}
\section*{3 · Musterlösung und Punkteraster}
\textbf{Allgemein:} Ganze Punkte je Zeile. Ein Fehler wird nur einmal abgezogen (Folgefehler). Gleichwertige Wege zählen voll.
In der Spalte P: \textbf{(E)}~= Ergebnispunkt, Ergebnis richtig \emph{und} Weg erkennbar; ohne Zeichen = die
Zeile bewertet den Weg, die Zeichnung oder die Begründung. Der ganze Test ist ohne Taschenrechner gestellt.
«Typische Fehler» nennt, welche Zeile 0 Punkte gibt, und die Punkte, die bleiben.
\textbf{Teilrichtige Zeilen:} Verlangt eine Zeile mehrere Angaben und stimmt nur ein Teil, gibt die Zeile 0 Punkte,
ausser die Zeile nennt ausdrücklich eine Aufteilung.
\textbf{Zahlen:} Dezimalzahl oder Bruch, wenn gleich ($4.5 = \tfrac{9}{2}$). Einheiten sind nicht verlangt.

\begin{lpaufgabe}{G1}{Werte}{3 P}
\begin{raster}
(a) Einsetzen & 1 & $-3 \cdot 2.5 + 12 = -7.5 + 12$\\
(a) Ergebnis & 1\E & $y = 4.5$\\
(b) & 1\E & $-3x + 12 = -6$, $-3x = -18$, $x = 6$\\
\end{raster}
\fehler{(a) Minus verloren ($3 \cdot 2.5 + 12 = 19.5$): beide Zeilen (a) 0 $\to$ höchstens 1 von 3\,P.
(b) $-6$ für $x$ eingesetzt ($y = 30$): Zeile (b) 0.
(b) $-3x = -18$ richtig, dann $x = -6$: Zeile (b) 0.}
\end{lpaufgabe}

\begin{lpaufgabe}{G2}{Akku}{3 P}
\begin{raster}
(a) Gleichung & 1 & $y = -6x + 80$ (Anfangswert $b = 80$, Änderung $m = -6$ pro Stunde)\\
(b) & 1 & $-6 \cdot 5 + 80 = 50$, also $50\,\%$. Auch als Folgepunkt mit der eigenen Gleichung, wenn richtig eingesetzt.\\
(c) & 1\E & $-6x + 80 = 20$, $-6x = -60$, $x = 10$ Stunden (auch ohne Gleichung: $80 - 20 = 60$, $60 : 6 = 10$)\\
\end{raster}
\fehler{$m$ und $b$ vertauscht ($y = 80x - 6$) oder Vorzeichen falsch ($y = 6x + 80$): Zeile (a) 0; (b) als Folgepunkt ($394$ bzw.\ $110$) 1 $\to$ höchstens 2 von 3\,P.
(c) $20 : 6$ gerechnet: Zeile (c) 0.}
\end{lpaufgabe}

\begin{lpaufgabe}{G3}{Steigung}{3 P}
\begin{raster}
Ansatz & 1 & $m = \dfrac{-3 - 5}{3 - (-1)}$: oben und unten dieselbe Reihenfolge ($B$ minus $A$ oder $A$ minus $B$)\\
Ergebnis & 1\E & $m = \dfrac{-8}{4} = -2$\\
Deutung & 1 & Die Gerade fällt (weil $m < 0$). Auch als Folgepunkt passend zum eigenen Vorzeichen.\\
\end{raster}
\fehler{Reihenfolge gemischt ($\tfrac{-3 - 5}{-1 - 3} = 2$): Ansatz- und Ergebniszeile 0; «steigt» als Folgepunkt 1 $\to$ 1 von 3\,P.
Bruch umgedreht ($\tfrac{4}{-8} = -\tfrac{1}{2}$): Ansatz- und Ergebniszeile 0; «fällt» 1 $\to$ 1 von 3\,P.
$3 - (-1) = 2$: Ergebniszeile 0 ($m = -4$) $\to$ 2 von 3\,P.}
\end{lpaufgabe}

\begin{lpaufgabe}{G4}{Zeichnen}{3 P}
\begin{center}
\gitter{-2}{9}{-4}{6}{\draw[lpblau,line width=1.2pt] (-2,-3.5) -- (9,4.75);
  \draw[lpgrau,line width=0.8pt] (0,-2) -- (4,-2) -- (4,1);
  \fill[lporange] (0,-2) circle (2.4pt); \fill (4,1) circle (2.2pt); \fill[lpgruen] (8,4) circle (2.4pt);}
\end{center}
\begin{raster}
Start & 1 & Punkt $(0 \mid -2)$ auf der $y$-Achse\\
Dreieck, Gerade & 1 & $4$ nach rechts, $3$ hinauf: $(4 \mid 1)$; Gerade durch beide Punkte (auch über Wertetabelle)\\
Punktprobe & 1\E & $\tfrac{3}{4} \cdot 8 - 2 = 6 - 2 = 4$: ja, $P$ liegt darauf\\
\end{raster}
\fehler{Start bei $(0 \mid 0)$ (Gerade $y = \tfrac{3}{4}x$): Startzeile 0; das Dreieck gibt seinen Punkt, wenn es $4$ nach rechts und $3$ hinauf geht $\to$ 2 von 3\,P.
Dreieck vertauscht ($3$ nach rechts, $4$ hinauf, zweiter Punkt $(3 \mid 2)$): Zeile «Dreieck, Gerade» 0.
Punktprobe nur abgelesen, ohne Rechnung: Zeile «Punktprobe» 0.}
\end{lpaufgabe}

\begin{lpaufgabe}{G5}{Ablesen}{3 P}
\begin{raster}
$b$ & 1 & $b = -2$ (Schnitt mit der $y$-Achse bei $(0 \mid -2)$)\\
$m$ & 1 & Von $(0 \mid -2)$ nach $(3 \mid 0)$: $3$ nach rechts, $2$ hinauf, $m = \tfrac{2}{3}$\\
Gleichung & 1\E & $y = \tfrac{2}{3}x - 2$\\
\end{raster}
\fehler{Bruch umgedreht ($m = \tfrac{3}{2}$): $m$-Zeile und Gleichungszeile 0 $\to$ 1 von 3\,P.
Nullstelle als $b$ genommen ($b = 3$): $b$-Zeile und Gleichungszeile 0 $\to$ 1 von 3\,P.}
\end{lpaufgabe}

\begin{lpaufgabe}{G6}{Infusion}{4 P}
\begin{center}
\sachgitter{6}{8}{50}{100}{\draw[lpblau,line width=1.2pt] (0,7.5) -- (5,0);
  \fill[lporange] (0,7.5) circle (2.4pt); \fill[lpgruen] (5,0) circle (2.4pt);}
\end{center}
\begin{raster}
(a) & 1 & $V = -3t + 750$\\
(b) & 1\E & $-3t + 750 = 0$, $t = 250$ Minuten\\
(c) Punkte & 1 & Start $(0 \mid 750)$ und Ende $(250 \mid 0)$ eingetragen; das Ende als Folgepunkt mit dem eigenen Wert aus (b)\\
(c) Strecke & 1 & gerade Strecke dazwischen, endet auf der $t$-Achse (nicht darunter weiter)\\
\end{raster}
\fehler{$V = 3t + 750$ oder $V = 750t - 3$: Zeile (a) 0; (b) gibt den Ergebnispunkt nur für $250$ mit erkennbarem Weg.
Gerade über $(250 \mid 0)$ hinaus ins Negative gezogen: Zeile «Strecke» 0 (ein Beutel enthält nie weniger als $0$ ml).}
\end{lpaufgabe}

\begin{lpaufgabe}{G7}{Zwei Punkte}{3 P}
\begin{raster}
$m$ & 1 & $m = \dfrac{7 - 1}{6 - 2} = \dfrac{6}{4} = \dfrac{3}{2}$\\
$b$ & 1 & $P$ eingesetzt: $1 = \tfrac{3}{2} \cdot 2 + b = 3 + b$, $b = -2$. Auch als Folgepunkt mit dem eigenen $m$.\\
Gleichung, Probe & 1\E & $y = \tfrac{3}{2}x - 2$; Probe mit $Q$: $\tfrac{3}{2} \cdot 6 - 2 = 9 - 2 = 7$. Beides verlangt.\\
\end{raster}
\fehler{$b$ mit einem $y$-Wert verwechselt ($b = 1$): $b$-Zeile und letzte Zeile 0 $\to$ 1 von 3\,P.
Bruch umgedreht ($m = \tfrac{2}{3}$, dann $b = -\tfrac{1}{3}$): $m$-Zeile 0, $b$-Zeile als Folgepunkt 1, letzte Zeile 0 $\to$ 1 von 3\,P.
Probe fehlt: letzte Zeile 0 $\to$ höchstens 2 von 3\,P.}
\end{lpaufgabe}

\begin{lpaufgabe}{G8}{Neue Rate}{3 P}
\begin{center}
\sachgitter{6}{7}{50}{100}{\draw[lpblau,line width=1.2pt] (0,6) -- (1,4) -- (5,0);
  \fill[lporange] (1,4) circle (2.4pt); \fill[lpgruen] (5,0) circle (2.4pt);}
\end{center}
\begin{raster}
(a) & 1 & $600 - 4 \cdot 50 = 400$ ml\\
(b) & 1\E & $400 : 2 = 200$ Minuten nach dem Umstellen, ab Beginn $50 + 200 = 250$ Minuten\\
(c) & 1 & Strecke $(0 \mid 600)$ bis $(50 \mid 400)$, dann flacher bis zum Ende aus (b): Knick bei $(50 \mid 400)$, kein Sprung. Das Ende als Folgepunkt mit dem eigenen Wert aus (b).\\
\end{raster}
\fehler{Nach dem Wechsel neu angefangen ($600 : 2 = 300$, leer nach $350$ Minuten, Graph springt auf $600$): Zeilen (b) und (c) 0 $\to$ 1 von 3\,P.
Wechsel übersehen ($600 : 4 = 150$ Minuten): Zeile (b) 0; (c) 0, wenn kein Knick gezeichnet ist.
Nur $200$ Minuten als Antwort (b), ohne die ersten $50$: Zeile (b) 0.}
\end{lpaufgabe}

\needspace{12\baselineskip}
\section*{4 · Selbsteinschätzung}
\begin{tabularx}{\linewidth}{@{}l X@{}}
\toprule
\textbf{22 – 25 P} & Die geprüften Teile sitzen. Wo du Punkte verloren hast: das Kapitel dieser Aufgabe nochmals.\\
\textbf{17 – 21 P} & Den schwächsten Teil nochmals, gemessen am Anteil (Teil A 6, B 9, C 10 P): Tüfteln und Übungen der Kapitel dieses Teils.\\
\textbf{11 – 16 P} & Zurück zu den Kapiteln aller Aufgaben, in denen du Punkte verloren hast.\\
\textbf{0 – 10 P} & Zurück zu den Kapiteln jedes Teils, in dem du weniger als die Hälfte erreicht hast; sind es alle drei, zum Vortest und zu Kapitel 1.\\
\bottomrule
\end{tabularx}

\smallskip
\textbf{Aufgabe $\to$ Kapitel:} G1, G2 $\to$ 1; G3 $\to$ 2; G4 $\to$ 3; G5 bis G8 $\to$ 4.
\end{document}
